%! AP Statistics — Unit 3, Lesson 3.5 — Guided Notes KEY
\documentclass[10pt]{article}
\usepackage{apstats-article}
\usepackage{apstats-key}

\begin{document}

\pageheader{Unit 3, Lesson 3.5}{Guided Notes — KEY}
\namedateperiod

\vspace{0.12in}

\begin{objectivebox}
\small By the end of this lesson, I will be able to\dots\\[4pt]
\ans{Identify the components of an experiment (units, factor, treatments, response, confounding).}\\[1pt]
\ans{Describe the four principles of a well-designed experiment.}\\[1pt]
\ans{Compare CRD, block design, and matched pairs design.}
\end{objectivebox}

\vspace{0.1in}

\begin{vocabbox}
\termblanklong{Experimental unit} \ans{The individual (person, animal, plot, etc.) assigned to a treatment.}
\termblanklong{Factor} \ans{The explanatory variable whose levels are intentionally manipulated by the researcher.}
\termblanklong{Treatment} \ans{A specific level (or combination of levels) of the factor applied to units.}
\termblanklong{Response variable} \ans{The outcome measured after the treatment is applied.}
\termblanklong{Control group / Placebo} \ans{A group receiving no active treatment (or an inactive look-alike) for comparison.}
\termblanklong{Random assignment} \ans{Assigning units to treatments by chance to balance confounding variables.}
\termblanklong{Blocking variable} \ans{A known source of variability; units are grouped by it so treatments are assigned within blocks.}
\termblanklong{Matched pairs design} \ans{A block design with blocks of size 2 — units matched on a key variable, or same unit gets both treatments.}
\end{vocabbox}

\vspace{0.1in}

\begin{notesbox}{Components of an Experiment (VAR-3.A)}
\small
\renewcommand{\arraystretch}{1.8}
\begin{tabularx}{\linewidth}{@{} >{\bfseries\color{navy}}p{0.3\linewidth} X @{}}
Experimental units & The \ans{individuals} assigned to treatments. \\
Factor & The \ans{explanatory} variable whose levels are \ans{manipulated} intentionally. \\
Treatment & A specific \ans{level} (or combination of levels) of the factor. \\
Response variable & The \ans{outcome} measured after treatments are applied. \\
Confounding variable & Related to both the \ans{factor} and the \ans{response}; may create a false association. \\
\end{tabularx}
\end{notesbox}

\vspace{0.1in}

\begin{notesbox}{Principles of a Well-Designed Experiment (VAR-3.B.1)}
\small
\begin{enumerate}[leftmargin=1.5em, itemsep=6pt]
  \item \textbf{Comparison:} at least \ans{two} treatment groups (one may be a \ans{control}).
  \item \textbf{Random assignment:} units assigned to treatments \ans{by chance}; this \ans{balances}
        confounding variables and allows \ans{causal} conclusions.
  \item \textbf{Replication:} more than \ans{one} unit per treatment group.
  \item \textbf{Control:} potential \ans{confounding} variables are managed (by design or by holding constant).
\end{enumerate}
\end{notesbox}

\vspace{0.1in}

\begin{notesbox}{Experimental Designs (VAR-3.C)}
\small
\renewcommand{\arraystretch}{1.9}
\begin{tabularx}{\linewidth}{@{} >{\bfseries\color{navy}}p{0.26\linewidth}
    >{\raggedright\arraybackslash}X @{}}
Completely randomized (CRD) &
All units randomly assigned to treatments --- \ans{no blocking}. \\
Block design &
Units grouped into \ans{homogeneous blocks} based on a blocking variable; treatments assigned
\ans{randomly} within each block. Reduces \ans{variability} due to the blocking variable. \\
Matched pairs &
A special block design; blocks of \ans{two} matched on a key variable (or same unit
gets \ans{both} treatments in random order). \\
Single-blind &
\ans{Subjects} don't know which treatment they received. \\
Double-blind &
Neither \ans{subjects} nor \ans{the evaluator/researcher} know treatment assignments. \\
\end{tabularx}
\end{notesbox}

\vspace{0.1in}

\begin{practicebox}
\small
\begin{enumerate}[leftmargin=1.5em, itemsep=6pt]
  \item Experimental units: \ans{60 college students}
  \item Factor: \ans{Drink type (energy drink vs.\ water)}
  \item Treatments: \ans{Energy drink; look-alike water (placebo)}
  \item Response variable: \ans{Reaction time}
  \item Type of blinding: \ans{Double-blind (students and researcher measuring don't know which drink)}
  \item Design type: \ans{Completely randomized design (CRD)}
\end{enumerate}
\end{practicebox}

\end{document}
